\documentclass[11pt]{article}
\usepackage[margin=25mm]{geometry}
\usepackage{fontspec}
\setmainfont{Latin Modern Roman}
\usepackage{amsmath,amssymb,mathtools}
\usepackage{unicode-math}
\setmathfont{Latin Modern Math}
\usepackage{xurl,hyperref}
\hypersetup{colorlinks=true,urlcolor=blue}
\setcounter{secnumdepth}{0}
\setlength{\parindent}{0pt}
\setlength{\parskip}{5pt}
\setlength{\emergencystretch}{3em}
\newcommand{\sref}[2]{\hyperlink{source:#1:#2}{[S#2]}}
\newcommand{\cataloguescope}[1]{\paragraph{Recorded target.}#1}
\newcommand{\R}{\mathbb R}
\newcommand{\C}{\mathbb C}
\newcommand{\N}{\mathbb N}
\newcommand{\Z}{\mathbb Z}
\newcommand{\Q}{\mathbb Q}
\newcommand{\F}{\mathbb F}
\newcommand{\eps}{\varepsilon}
\begin{document}
\section{Canonical versus abstract hyper-Kähler structures on cotangent bundles}\label{statement}

\subsection{Definitions and mathematical statement}

\paragraph{Shared notation.}$\mathbb N=\{1,2,\ldots\}$; the standard number fields have their usual meanings. All additional conventions are stated locally.


\paragraph{Definitions and notation.}A Kähler manifold has a complex structure $I$ and a compatible Riemannian metric whose associated two-form is closed. A hyper-Kähler structure consists of a Riemannian metric and three parallel compatible complex structures $I,J,K$ satisfying $I^2=J^2=K^2=IJK=-1$. The holomorphic cotangent bundle $T^*M$ also has a canonical holomorphic symplectic form. Compatibility with that canonical structure, completeness and existence merely near the zero section are additional, distinct conditions.

\paragraph{Catalogue wording retained for traceability.}
Hyper-K\"{a}hler Structure on $T^* M$

Let $M$ be a K\"{a}hler manifold. Must $T^*M$ admit a hyper-K\"{a}hler structure?
\paragraph{Scope and source qualification.}Primary AIM Problem 2.1 explicitly reports that counterexamples have been found. The inherited partially_solved label is retained as catalogue metadata, without claiming that this universal existence question remains open. The source does not specify canonical compatibility or completeness; a local construction does not answer global existence. The catalogue title/status was paired with a machine-generated research attempt; that attempt is not peer-reviewed evidence of settlement. The statement covers the original source questions rather than only the title’s subcase.
\subsection{Short English statement}
Determine whether the entire cotangent bundle of every Kähler manifold admits a hyper-Kähler structure.
\subsection{Research attempt}
\paragraph{Model, date and investigated scope.}
OpenAI GPT-6 Astra (configured model \texttt{gpt-6-astra}), ultra reasoning; generated 11 October 2026. Constructed a complete compatible flat hyper-Kahler metric on cotangent bundles of flat complex tori. The source-reported universal counterexamples are preserved, not newly verified or contradicted.

\paragraph{Sources and literature check.}
Read the supplied complete definition, catalogue wording, scope qualifications and source list for this canonical ID. The original-source citations are inherited as supplied references; no fresh primary-source or novelty search was performed for this attempt. No current open-status claim is inferred from this calculation.

\paragraph{A global flat positive example.}Let \(M=\mathbb C^m/\Lambda\), where \(\Lambda\) is any full lattice of translations, equipped with the quotient flat Kahler metric. Its holomorphic cotangent bundle is \((\mathbb C^m_z\times\mathbb C^m_w)/\Lambda\), where \(\lambda\) acts by \((z,w)\mapsto(z+\lambda,w)\). On the cover take the Euclidean metric and real-linear maps
\[I(z,w)=(iz,iw),\qquad J(z,w)=(-\overline w,\overline z),\qquad K=IJ.
\]
Their derivatives are constant, \(I^2=J^2=-I_{\mathrm{id}}\), and \(IJ=-JI\). They are orthogonal and parallel, so their metric two-forms are closed. All these tensors are invariant under translations and descend to the cotangent bundle, giving a complete hyper-Kahler metric. Completeness follows because the covering Euclidean metric is complete and the translations act by isometries.
\paragraph{Compatibility check.}In these coordinates the canonical holomorphic symplectic form is \(\sum_j dz_j\wedge dw_j\), up to the conventional overall sign for the tautological form. The Euclidean hyper-Kahler structure produces this form up to a constant phase/sign choice of \(J,K\), so this particular construction is compatible. It uses the affine transition functions critically: for general base coordinate changes, the conjugate-linear exchange defining \(J\) does not descend by the same argument.
\paragraph{Scope.}The supplied primary-source qualification already reports counterexamples to the universal existence question. This flat construction neither disproves that report nor checks the counterexamples. It supplies a familiar global positive subcase and shows exactly which transition-function feature makes it work. Whether a structure is canonical, complete, local near the zero section or merely abstract remains a separate question; none is silently substituted for the full global question in the source.

\paragraph{Verification and adversarial audit.}
Checked quaternionic relations, parallel orthogonal tensors, translation descent, completeness and canonical symplectic compatibility up to sign. This was an LLM adversarial audit, not independent human verification; no theorem prover was used. Prompts A then B applied to this ID: Checked quaternionic relations, parallel orthogonal tensors, translation descent, completeness and canonical symplectic compatibility up to sign. Scope remains limited: Analyze non-affine transition functions or inspect the reported counterexamples; the flat construction does not extend automatically.

Coordinator checked the flat quotient tensor construction, quaternionic identities, completeness and compatibility with the canonical holomorphic symplectic form. This additional check was model-assisted, not independent human verification.

\paragraph{Tools, process and contributions.}
Python 3 and local shell used to read supplied files and serialize this manuscript; mathematics derived symbolically in the text, with no numerical experiment or formal proof assistant. Contribution prompt A was followed by prompt B. GPT-6 Astra with ultra reasoning drafted the research and performed its stated self-audit; any additional coordinator contribution is identified in the verification and audit. A GPT-6-based Codex coordinator assembled the manuscript and checked repository integration; its exact variant and reasoning configuration were not exposed. No human mathematical contribution, independent human proof check or formal proof certification is supplied. Complete derivations and the consulted source references are included here; private reasoning transcripts are not published.

\paragraph{Final status and first remaining gap.}
\textbf{UNRESOLVED}. Analyze non-affine transition functions or inspect the reported counterexamples; the flat construction does not extend automatically.

\section{Completion Estimate}
\noindent\textbf{COMPLETION: 0\%}
The calculation is an elementary benchmark or conditional reduction; it establishes no claimed new portion of the complete source target.

\subsection{Sources}
\begin{itemize}
\item[\hypertarget{source:statement:1}{S1}] ULAM.AI and the UnsolvedMath Contributors, UnsolvedMath, catalogue entry. Catalogue curation under CC BY 4.0; definitions and English exposition added locally.
\url{https://huggingface.co/datasets/ulamai/UnsolvedMath}.
\item[\hypertarget{source:statement:2}{S2}] American Institute of Mathematics, Non-Hermitian quantum mechanics and symplectic geometry, source item 2.1.
\url{http://aimpl.org/nhquantumsymp/2/}.
\end{itemize}
\label{end:statement}
\end{document}
